\SourceLeaf{414}%
aus \(-\frac12b[2]+\frac12a[3]\) und \(\frac12B[2]-\frac12A[3]\) folgt, wenn man zur Abkürzung
\[\begin{gathered}aB-bA=\lambda,\quad ab(A+B)=f,\quad\tfrac12AB(A+B)+\lambda n=g\\
AB(a+b)=F,\quad\tfrac12aBB+\tfrac12bAA=G\end{gathered}\]
schreibt
\[\begin{gathered}\lambda T+G=\lambda(t+f)x,\quad\lambda t-g=(\lambda T-F)x\quad\text{oder}\quad\lambda(t-Tx)=g-Fx\end{gathered}\]
und folglich
\[\begin{aligned}\lambda t&=\frac{+g-(F+G)x+fxx}{1-xx}\quad.\ .\ .\ .\ .\ .\ .\ .\ [4]\\
\lambda T&=\frac{-G+(f+g)x-Fxx}{1-xx}\quad.\ .\ .\ .\ .\ .\ .\ .\ [5]\end{aligned}\]
\par
Die Gleichung \([1]\) in der Form \((t-Tx)^2+TT(1-xx)=rr\) gibt
\[\lambda\lambda rr-(g-Fx)^2=\lambda\lambda TT(1-xx)\quad.\ .\ .\ .\ .\ .\ [6]\]
\par
Substituirt man darin den Werth von \(\lambda T\) aus \([5]\), so erhält man die cubische Gleichung
\[\begin{aligned}&2fFx^3-(ff+2fg+FF+2FG+\lambda\lambda rr)xx\\
&\quad+(2fG+2gF+2gG)x+\lambda\lambda rr-gg-GG=0\quad.\ .\ [7]\end{aligned}\]
\par
Nachdem dadurch \(x\) bestimmt ist, findet man die Coordinaten des Punktes \((2)\) aus obigen Ausdrücken, die, wenn man darin für \(t-Tx\) und \(T\) die vorhin gegebenen Werthe substituirt, folgende Gestalt annehmen
\[n-\frac{g-Fx}{\lambda},\qquad\frac{-G+(f+g)x-Fxx}{\lambda\BareRoot{(1-xx)}}\]
und die beiden andern Punkte \((1)\) und \((3)\), indem man zu diesen Werthen \(-Ax,-A\BareRoot{(1-xx)}\) und \(+Bx,+B\BareRoot{(1-xx)}\) beziehungsweise hinzufügt.

Da jedem Cosinus zwei Werthe des Winkels angehören, oder was dasselbe ist, da die Radical-Grösse \(\BareRoot{(1-xx)}\) sowol positiv wie negativ genommen werden kann, so gibt jede zulässige Wurzel der Gleichung zwei Auflösungen; nemlich zwei gegen die Linie \(\mathrm{(II)}\) symmetrische Lagen der Punkte \((1),(2),(3)\), was auch schon für sich klar ist. Für den Fall, dass \(+1\) oder \(-1\) eine Wurzel der Gleichung \([7]\) wird, ist übrigens die obige Formel für die Ordinate nicht brauchbar, weil dann Zähler und Nenner null werden, und man muss anstatt
